\section{Unequal twists through beta mixtures and finite resolvent identities}
\label{mxtw:sec:main}
The fixed-twist calculus suggests a next question: how do kernels with
different frequency shifts convolve? The Fourier multipliers give an
explicit answer. Fractional orders admit a Dirichlet beta mixture of
intermediate twists, while positive integer orders reduce to finitely
many elementary kernels. These statements use the classical beta
integral and partial fractions; they make no literature-priority or
arithmetic-independence claim.

Write $K_\alpha^\theta=W_{1-\alpha}^\theta$, so that
\[
 \widehat K_\alpha^\theta(n)
       =(2\pi i(n+\theta))^{-\alpha}.
\]
The twists in the fractional-mixture theorem lie in $(0,1)$, with the
logarithms of \eqref{tw:eq:lambda}. More generally one may put them in a common
open interval between consecutive integers. This condition keeps the
frequencies at each fixed mode on the same imaginary ray.

\begin{theorem}[Any finite number of fractional twist kernels]
\label{mxtw:thm:mixture}
Let $d\ge2$, $\theta_j\in(0,1)$ and
$\operatorname{Re}\alpha_j>0$. Put $A=\sum_j\alpha_j$ and let
$\Sigma_{d-1}=\{t_j\ge0:\sum_jt_j=1\}$, with measure
$dt_1\cdots dt_{d-1}$. Then, as periodic distributions,
\begin{equation}\label{mxtw:eq:mixture}
 K_{\alpha_1}^{\theta_1}*\cdots*K_{\alpha_d}^{\theta_d}
 =\frac{\Gamma(A)}{\prod_j\Gamma(\alpha_j)}
   \int_{\Sigma_{d-1}}
       \left(\prod_jt_j^{\alpha_j-1}\right)
       K_A^{\sum_jt_j\theta_j}\,dt.
\end{equation}
The mixture includes the zero Fourier coefficient. Its normalized
right side, continued as a whole, is entire in all $\alpha_j$.
For two twists the simplex is the interval $0\le t\le1$ and the
intermediate twist is $t\theta_1+(1-t)\theta_2$.
\end{theorem}
\begin{proof}
For positive real $a_j$ and $\operatorname{Re}\alpha_j>0$, change
variables $y_j=rt_j$ in the product of Gamma integrals
\[
 \prod_j a_j^{-\alpha_j}
 =\frac1{\prod_j\Gamma(\alpha_j)}
   \int_{(0,\infty)^d}\prod_j y_j^{\alpha_j-1}
                         e^{-\sum_j a_jy_j}\,dy.
\]
The Jacobian is $r^{d-1}$. Integrating $r$ gives the scalar
Dirichlet formula with $(\sum_jt_ja_j)^{-A}$.
At Fourier mode $n$, all $n+\theta_j$ have the same sign.
Apply the scalar formula to $a_j=|2\pi(n+\theta_j)|$ and restore
their common phase $\exp(-i\pi A\operatorname{sgn}(n+\theta_j)/2)$.
The resulting coefficient is exactly that of
\eqref{mxtw:eq:mixture}, including $n=0$.

The intermediate twists range in a compact subinterval of $(0,1)$.
On compact parameter subsets, the kernel coefficients have a common
polynomial growth bound, and the simplex weight is integrable.
Rapid decrease of the Fourier coefficients of each smooth test
function therefore justifies integration and summation. Fourier
uniqueness proves the distribution identity. For continuation outside
the initial half-planes, subtract finite Taylor polynomials at the
simplex faces before integrating; the added monomials have meromorphic
beta integrals. Their combined normalized continuation has the same
Fourier coefficients as the left side, which are entire in every
$\alpha_j$. Test-function pairing gives normal convergence on compact
parameter sets. Thus all apparent poles of the combined continuation
are removable. It is this continuation, not a divergent unsubtracted
simplex integral, that is asserted to be entire.
\end{proof}

\begin{theorem}[Finite reduction at every positive integer order]
\label{mxtw:thm:integer}
For distinct real nonintegral $\theta,\phi$, set
$\Delta=2\pi i(\phi-\theta)$. For positive integers $r,s$,
\begin{align}\label{mxtw:eq:partial-fractions}
 K_r^\theta*K_s^\phi
 ={}&\sum_{j=1}^r
   \frac{(-1)^{r-j}\binom{r+s-j-1}{r-j}}
        {\Delta^{r+s-j}}K_j^\theta\notag\\
 &+\sum_{k=1}^s
   \frac{(-1)^r\binom{r+s-k-1}{s-k}}
        {\Delta^{r+s-k}}K_k^\phi.
\end{align}
Every ordinary representative on $(0,1)$ is elementary, since
\begin{equation}\label{mxtw:eq:elementary}
 K_1^\theta(x)=\frac{e^{-2\pi i\theta x}}{1-e^{-2\pi i\theta}},
 \qquad
 K_r^\theta(x)=
 \frac{(-1)^{r-1}}{(2\pi i)^{r-1}(r-1)!}
           \partial_\theta^{r-1}K_1^\theta(x).
\end{equation}
If two twists coincide, their orders add:
$K_r^\theta*K_s^\theta=K_{r+s}^\theta$.
\end{theorem}
\begin{proof}
At a Fourier mode put $X=2\pi i(n+\theta)$; the second multiplier
has denominator $X+\Delta$. The principal part at $X=0$ of
$X^{-r}(X+\Delta)^{-s}$ is obtained by expanding
$(X+\Delta)^{-s}$ through degree $r-1$; this gives the first sum.
At $X=-\Delta$, expand $X^{-r}$ in $X+\Delta$ through degree
$s-1$; this gives the second sum. Their difference from the rational
function has no finite poles and tends to zero at infinity, so it is
zero. Apply this identity at every Fourier mode and use uniqueness.
The representative for $K_1$ is \eqref{tw:eq:Lerch} at $s=0$.
Differentiation of its multiplier gives
$\partial_\theta K_r^\theta=-2\pi i rK_{r+1}^\theta$;
distributional differentiation in the parameter is legitimate by the
same compact Fourier estimates. This proves
\eqref{mxtw:eq:elementary}. The coincident identity is
\eqref{tw:eq:group}.
\end{proof}

For instance, the first unequal-twist convolution is the explicit
ordinary integral
\[
 (K_1^\theta*K_1^\phi)(x)
   =\frac{K_1^\theta(x)-K_1^\phi(x)}{2\pi i(\phi-\theta)}.
\]
Its coincident limit is
$-\partial_\theta K_1^\theta/(2\pi i)=K_2^\theta$.
Any finite number of distinct twists and positive integer orders is
handled by the same principal-part argument; repeated twists are
combined first. These integer-order reductions require no common
frequency interval. Fractional orders have the precise mixture
\eqref{mxtw:eq:mixture}. Finite reduction of every mixed Stieltjes
order jet to a prescribed list of single Hurwitz values remains a
separate question; the mixture does not assert such a reduction.
